\documentclass[12pt]{article}
\usepackage[margin=0.75in]{geometry}
\usepackage{fontspec}
\setmainfont{Latin Modern Roman}
\usepackage{microtype}
\usepackage{multicol}
\usepackage{fancyhdr}
\usepackage{titlesec}
\usepackage{enumitem}
\usepackage{xcolor}
\usepackage[hidelinks]{hyperref}
\setlength{\columnsep}{0.28in}
\setlength{\parindent}{1.1em}
\setlength{\parskip}{0.32em}
\setlength{\headheight}{15pt}
\titleformat{\section}{\large\bfseries\scshape}{}{0pt}{}
\titleformat{\subsection}{\normalsize\bfseries}{}{0pt}{}
\pagestyle{fancy}
\fancyhf{}
\fancyfoot[C]{\thepage}
\renewcommand{\headrulewidth}{0.4pt}
\newcommand{\focusbox}[1]{\par\vspace{0.4em}\noindent\begingroup\setlength{\fboxsep}{6pt}\colorbox{gray!12}{\begin{minipage}{\dimexpr\linewidth-2\fboxsep\relax}#1\end{minipage}}\endgroup\par\vspace{0.5em}}

\fancyhead[L]{Whately: Discovery and Judgment}
\fancyhead[R]{Student Reading}
\begin{document}
\begin{center}
{\LARGE\bfseries Reasoning, Discovery, and Judgment}\par\vspace{0.25em}
{\large\scshape Week 14: Can Logic Discover Anything New?}\par\vspace{0.5em}
{Adapted and modernized from Richard Whately, \textit{Elements of Logic}}\par\vspace{0.6em}
\rule{0.82\textwidth}{0.4pt}
\end{center}

\section*{Central Question}
Does reasoning discover truth, or does it only arrange what we already know?

\section*{Vocabulary Preview}
\begin{multicols}{2}
\begin{itemize}[leftmargin=*, itemsep=0.18em]
\item \textbf{Discovery}: finding or learning something not yet known.
\item \textbf{Proof}: showing that a conclusion follows from given reasons.
\item \textbf{Observation}: direct attention to facts, events, experiments, or evidence.
\item \textbf{Testimony}: information received from another source or witness.
\item \textbf{Memory}: retained knowledge that can supply material for reasoning.
\item \textbf{Imagination}: the mind's power to propose possibilities, patterns, or hypotheses.
\item \textbf{Judgment}: deciding what is relevant, probable, or wise in a particular case.
\item \textbf{Inference}: moving from premises to a conclusion.
\item \textbf{Method}: an ordered way of searching, testing, and arranging evidence.
\item \textbf{Intellectual-tool map}: a map that names which mental work is needed: observe, define, infer, test, or judge.
\end{itemize}
\end{multicols}

\section*{Introduction}
By Week 14, you have used many tools: term-definition, proposition translation, syllogism mapping, validity tests, fallacy autopsies, analogy tests, induction checks, and credibility maps. It would be easy to exaggerate logic and treat it as a machine that discovers every truth by itself.

Whately warns against that overclaim. Logic is powerful, but it has a proper office. It tests and arranges reasoning. It helps us see whether a conclusion follows from what has been stated. But it does not replace observation, testimony, memory, imagination, experience, or wise judgment.

\focusbox{\textbf{Source note:} Adapted and modernized from Whately's discussion of the discovery of truth, the province of logic, and the limits of reasoning alone. Francis Bacon's emphasis on observation and method stands behind several classroom examples, but the handout remains Whately-centered. Classroom examples are original.}

\begin{multicols}{2}
\section*{Reading}

\subsection*{1. Logic Has a Proper Office}
Logic does not create facts out of nothing. It asks whether one thought follows from another. If the premises are unclear, logic asks for clearer terms. If the argument has a bad form, logic exposes the defect. If the conclusion is larger than the evidence, logic asks for a narrower claim.

This is not a small task. Much bad thinking survives because people skip these tests. But the task is not the whole of knowledge.

\subsection*{2. Reasoning Needs Material}
An argument needs material to reason from. That material may come from observation, testimony, memory, records, experiments, reading, or ordinary experience.

If a class wants to know whether the courtyard floods after heavy rain, a beautiful syllogism is not enough. Someone must look at the courtyard, inspect records, ask witnesses, or measure the problem. Reasoning can organize the evidence, but it cannot replace the evidence.

\subsection*{3. Discovery and Proof Are Not the Same}
Discovery answers, ``How did we find or learn this?'' Proof answers, ``Why should this conclusion be accepted?'' The two often work together, but they are not identical.

A student may discover a pattern by noticing that several successful speeches begin with a concrete story. The student still has to prove any larger claim about speeches. A scientist may imagine a hypothesis before proving it. A historian may discover a letter in an archive, then reason about what it shows.

\subsection*{4. Observation and Bacon's Warning}
Francis Bacon warned that the mind often leaps too quickly from a few examples to a broad theory. Whately's logic gives a similar caution in another form: do not let the conclusion outrun the evidence.

Observation matters because it gives the mind something outside itself to answer to. Method matters because observation can be careless, selective, or biased. Good inquiry asks: What have we actually seen? What have we not seen? What would count against our first guess?

\subsection*{5. Testimony, Memory, and Imagination}
Not every discovery begins with direct observation. You may learn from a witness, a book, a map, a record, or a teacher. You may remember an earlier case. You may imagine a possible explanation and then test it.

Imagination is not proof, but it is not useless. It can propose a hypothesis. Memory is not proof, but it can supply relevant cases. Testimony is not proof by itself, but credible testimony may give strong evidence. Logic helps sort these materials.

\subsection*{6. Judgment Applies the Tool}
Judgment decides which tool is needed. A verbal dispute needs definition. A rumor needs source-checking. A broad claim from a few examples needs an induction test. A formal syllogism needs a validity check. A literary interpretation needs evidence from the text and a claim sized to that evidence.

A disciplined thinker does not use one tool for every problem. The mature question is not, ``Can I force this into a syllogism?'' The mature question is, ``What kind of intellectual work does this problem need?''

\subsection*{7. The Intellectual-Tool Map}
Use this Week 14 map:
\begin{enumerate}[leftmargin=*,itemsep=0.12em]
\item What is the question or problem?
\item What must be discovered by observation, testimony, memory, or research?
\item What terms or claims must be clarified?
\item What inference or proof must be tested?
\item What judgment is needed about relevance, probability, or application?
\end{enumerate}

\subsection*{8. Literary Cases}
In \textit{Macbeth}, characters discover too late that the witches' words did not mean what Macbeth wanted them to mean. In \textit{Julius Caesar}, Rome discovers too late how dangerous it is when political judgment is driven by fear, honor, signs, and speeches before the evidence has been weighed carefully.

The tragedy is not that characters never reason. They do reason. The tragedy is that they use the wrong tool, stop too soon, trust the wrong sign, ignore needed evidence, or judge before discovery is complete.

\section*{Reading Focus}
\begin{enumerate}[leftmargin=*]
\item What is logic's proper office?
\item Why does reasoning need material from outside itself?
\item What is the difference between discovery and proof?
\item Why does observation need method?
\item How can imagination help inquiry without becoming proof?
\item What does judgment decide when several logic tools are available?
\item Copy the five steps of the intellectual-tool map.
\end{enumerate}

\section*{Handout Summary}
Write 5--7 sentences explaining why logic is necessary but not sufficient for discovering truth. Your summary must include the words \textit{discovery}, \textit{proof}, \textit{observation}, and \textit{judgment}.
\end{multicols}
\end{document}
